\chapter{熱的 Green 関数}
\label{chap:thermal-green}
\chapterepigraph{``The energy of the universe is constant. The entropy of the universe strives toward a maximum.''}{Rudolf Clausius, ``Concerning Several Conveniently Applicable Forms for the Main Equations of the Mechanical Heat Theory'' (1865)}

\section{KMS 条件}

時間発展を生成する Hamiltonian を $H$ とし、\term{逆温度} $\beta$ の\term{平衡状態}を
\begin{equation}
\rho_\beta
=
\frac{\rme^{-\beta H}}{Z},
\qquad
Z
=
\Tr\rme^{-\beta H}
\label{eq:gibbs-state}
\end{equation}
とする。\term{Hermite 演算子} $O(t)$ に対して
\begin{align}
G^>(t)
&=
\Tr\left(
\rho_\beta O(t)O(0)
\right),
\\
G^<(t)
&=
\Tr\left(
\rho_\beta O(0)O(t)
\right)
\end{align}
を定義する。\term{trace} の巡回性と\term{虚時間発展}から
\begin{equation}
G^>(t)
=
G^<(t+\rmi\beta)
\label{eq:kms-time}
\end{equation}
を得る。これを \term{Kubo--Martin--Schwinger（KMS）条件}とよぶ。有限系の Gibbs 行列を
書けない場の理論でも、\term{相関関数}の解析性と式\eqref{eq:kms-time}を\term{熱平衡}の定義に
できる。

Fourier 変換を
\begin{equation}
G(\omega)
=
\int_{-\infty}^{\infty}\rmd t\,
\rme^{+\rmi\omega t}G(t)
\label{eq:thermal-fourier}
\end{equation}
とすれば、\term{KMS 条件}は \term{detailed balance}
\begin{equation}
G^>(\omega)
=
\rme^{\beta\omega}G^<(\omega)
\label{eq:kms-frequency}
\end{equation}
となる。

\section{\term{Lehmann 表示}}

$H\ket{n}=E_n\ket{n}$ とし、$O_{mn}=\bra{m}O\ket{n}$ とする。\term{完全系}を二点関数へ
挿入すれば、
\begin{align}
G^>(\omega)
&=
\frac{2\pi}{Z}
\sum_{m,n}
\rme^{-\beta E_n}
|O_{mn}|^2
\delta\bigl(\omega-E_m+E_n\bigr),
\label{eq:lehmann-greater}
\\
G^<(\omega)
&=
\frac{2\pi}{Z}
\sum_{m,n}
\rme^{-\beta E_m}
|O_{mn}|^2
\delta\bigl(\omega-E_m+E_n\bigr)
\label{eq:lehmann-less}
\end{align}
を得る。delta 関数の上では $E_m-E_n=\omega$ なので、項ごとに
$G^>(\omega)=\rme^{\beta\omega}G^<(\omega)$ が成り立つ。\term{有限体積}ではスペクトルは
\term{離散線}の和であり、\term{熱力学極限}または有限分解能の\term{粗視化}によって滑らかな
\term{スペクトル関数}になる。

この表示は、熱的 Green 関数が抽象的な解析関数であるだけでなく、観測量の行列要素と
\term{状態密度}の和であることを示す。次章の ETH は、式\eqref{eq:lehmann-greater}に現れる
個々の $O_{mn}$ に滑らかな包絡を仮定し、一つの高エネルギー固有状態から同じ\term{詳細
釣り合い}が現れる条件を問う。

\section{スペクトル関数と揺らぎ}

\term{スペクトル関数}と\term{対称相関関数}を
\begin{align}
\rho_O(\omega)
&=
G^>(\omega)-G^<(\omega),
\label{eq:thermal-spectral-density}
\\
F_O(\omega)
&=
\frac{1}{2}
\left[
G^>(\omega)+G^<(\omega)
\right]
\label{eq:thermal-symmetric-density}
\end{align}
と定義する。式\eqref{eq:kms-frequency}から
\begin{equation}
F_O(\omega)
=
\frac{1}{2}
\coth\left(\frac{\beta\omega}{2}\right)
\rho_O(\omega)
\label{eq:fdt-rho}
\end{equation}
を得る。

本書の遅延 Green 関数は
\begin{equation}
G_R(t)
=
\rmi\theta(t)
\Tr\left(
\rho_\beta[O(t),O(0)]
\right)
\end{equation}
である。実周波数の境界値に対し
\begin{equation}
\rho_O(\omega)
=
2\im G_R(\omega)
\label{eq:spectral-im-retarded}
\end{equation}
なので、\term{揺動散逸関係}は
\begin{equation}
F_O(\omega)
=
\coth\left(\frac{\beta\omega}{2}\right)
\im G_R(\omega)
\label{eq:fluctuation-dissipation}
\end{equation}
となる。応答の\term{散逸部分}が、平衡状態の量子・熱揺らぎを定める。

\begin{warningbox}
式\eqref{eq:spectral-im-retarded}の符号は遅延関数と Fourier 変換の規約に依存する。
文献の式を移すときは、$G_R$ の前の $\rmi$、指数 $\rme^{\pm\rmi\omega t}$、
\term{スペクトル関数}の順序を一組として確認する。
\end{warningbox}

\section{吸収と放出}

外部の弱い周期源 $j(t)$ を $-j(t)O(t)$ と結合する。単位時間当たりの吸収は、
$\omega>0$ に対して $\rho_O(\omega)$、すなわち $\im G_R(\omega)$ に比例する。
一方、状態自身の\term{自発・誘導放出}は $G^<(\omega)$ によって測られる。\term{KMS 条件}は
\begin{equation}
\frac{G^<(\omega)}{G^>(\omega)}
=
\rme^{-\beta\omega}
\label{eq:detailed-balance-ratio}
\end{equation}
を要求する。

ブラックホールでは、古典的\term{吸収断面積}または greybody factor が $\rho_O$ の実軸上の
散逸を与える。Hartle--Hawking 状態では \term{KMS 条件}が吸収と放出を結び、Unruh 状態では
外向き sector に \term{Hawking 占有数}が現れるが、大域的な\term{熱平衡}とはならない。

\begin{principlebox}
古典吸収、量子放出、\term{平衡ノイズ}は三つの独立な表ではない。遅延 Green 関数の虚部と
KMS 条件が、それらを同じスペクトル関数へまとめる。
\end{principlebox}

\section{QNM と熱極を混同しない}

有限温度計算では \term{Euclid 周波数}
\begin{equation}
\omega_n^{E}
=
\frac{2\pi n}{\beta},
\qquad
n\in\bbZ
\label{eq:bosonic-matsubara}
\end{equation}
が現れる。これは虚時間の周期性から生じる \term{Matsubara 周波数}であって、放射境界条件の
QNM 周波数ではない。\term{Euclid 相関関数}を適切に解析接続すれば遅延関数を得て、その
別の sheet の極として QNM が現れる。

したがって、\term{Matsubara 点}の列と QNM 極の列が図の上で似ていても同一視してはならない。
前者は熱状態のサンプリング点、後者は実時間応答の共鳴である。両者を結ぶには解析
接続の一意性と高周波挙動が必要になる。

\section{回転する熱平衡}

定常回転系では、密度演算子は模式的に
\begin{equation}
\rho
\mathrel{\propto}
\exp\left[
-\beta
\left(
H-\Omega_H J-\Phi_H Q
\right)
\right]
\label{eq:rotating-gibbs}
\end{equation}
となる。KMS 因子には式\eqref{eq:horizon-effective-frequency}の
$\widetilde\omega$ が入る。漸近平坦 Kerr では \term{superradiant mode} のため、全外部領域を
覆う通常の Hartle--Hawking 状態の構成には障害がある。熱因子を局所的に用いることと、
大域的な \term{Gibbs 状態}の存在を区別する。

\section{本章の結論}

熱状態は KMS 条件で特徴づけられ、\term{揺動散逸関係}が状態依存の揺らぎを状態に依存しない
遅延応答へ結ぶ。ブラックホールの古典吸収と量子放出は、この関係の両側に位置する。
次章では、熱平衡を一つの固有状態の行列要素から説明しようとする ETH を導入する。
